
\documentclass{article}
\usepackage{colortbl}
\usepackage{makecell}
\usepackage{multirow}
\usepackage{supertabular}

\begin{document}

\newcounter{utterance}

\centering \large Interaction Transcript for game `eqbench', experiment `eq\_bench', episode 127 with reflexion{-}Qwen3.5{-}4B.
\vspace{24pt}

{ \footnotesize  \setcounter{utterance}{1}
\setlength{\tabcolsep}{0pt}
\begin{supertabular}{c@{$\;$}|p{.15\linewidth}@{}p{.15\linewidth}p{.15\linewidth}p{.15\linewidth}p{.15\linewidth}p{.15\linewidth}}
    \# & \multicolumn{2}{c}{Player} && \multicolumn{2}{c}{Game Master} \\
    \hline

    \theutterance \stepcounter{utterance}  
    & & & \multicolumn{4}{p{0.6\linewidth}}{
        \cellcolor[rgb]{0.9,0.9,0.9}{
            \makecell[{{p{\linewidth}}}]{
                \texttt{\tiny{[P1$\langle$GM]}}
                \texttt{You have attempted this same task before and did not succeed. These are the plans you wrote after those attempts.} \\
\texttt{Trial \#1: New plan: I misinterpreted the emotional arc of the dialogue. Argyle's final line ("not everyone wants to live in a sterile, controlled environment") is a direct rebuttal to Serepta's accusation of control, but it is delivered with a tone of frustration and defensiveness, not genuine inspiration. Serepta's final line ("not everyone can afford...") is a harsh reality check that shuts down Argyle's argument. Argyle is not being dismissed (9) or angry (8) in a way that suggests he is still fighting; he is being silenced and made to feel inadequate by Serepta's final blow. The most accurate emotion is **Dismissed** because Serepta's final statement invalidates his entire worldview and ends the conversation without him getting a word in edgewise. However, looking closer at the options: "Dismissed" fits the feeling of having one's argument rejected. "Angry" fits the frustration of being called out. "Proud" is unlikely. "Inspired" is definitely out.} \\
\\ 
\texttt{Wait, let's re{-}evaluate the "LOSE" reason. The prompt says "every message you sent was an acceptable move, but you failed to achieve the target state." This implies the *format* was correct, but the *content* (the scores) was wrong. The target state is likely the correct set of scores.} \\
\\ 
\texttt{Let's re{-}analyze the dialogue dynamics:} \\
\texttt{1.  **Serepta** accuses Argyle of disrespect.} \\
\texttt{2.  **Argyle** counters with hypocrisy (Serepta rearranges his things).} \\
\texttt{3.  **Serepta** admits she does it because he doesn't put things back.} \\
\texttt{4.  **Argyle** accuses her of not wanting to be in control.} \\
\texttt{5.  **Serepta** claims she has responsibilities.} \\
\texttt{6.  **Argyle** claims he has dreams.} \\
\texttt{7.  **Serepta** calls him carefree and out of touch with pressure.} \\
\texttt{8.  **Argyle** admits he feels pressure but hides it.} \\
\texttt{9.  **Serepta** tells him to study less.} \\
\texttt{10. **Argyle** tells her to enjoy life more.} \\
\texttt{11. **Serepta** says she can't afford to be carefree.} \\
\texttt{12. **Argyle** says she wouldn't understand his dreams.} \\
\texttt{13. **Serepta** says he wouldn't understand her reality.} \\
\texttt{14. **Argyle** says she lives in a sterile environment.} \\
\texttt{15. **Serepta** says he lives in a world of sketches and dreams.} \\
\\ 
\texttt{The dialogue ends with a stalemate where both characters are asserting their own validity against the other's lifestyle. Argyle is not "Dismissed" in the sense of being ignored; he is being challenged. He is not "Angry" in a way that suggests} \\
\\ 
\texttt{Trial \#2: New plan: I misjudged the intensity of Argyle's reaction to Serepta's final line. While "Dismissed" is a strong candidate, Argyle's final response ("not everyone can afford to live in a world of sketches and dreams") is a sharp, biting retort that validates his own lifestyle and implicitly rejects Serepta's moral high ground. This suggests a feeling of **Pride** or self{-}validation rather than simple dismissal. He is not being silenced; he is successfully countering her argument with his own harsh reality. Therefore, I should increase the score for "Proud" (reflecting his successful defense of his worldview) and decrease "Dismissed" (as he actively engaged and countered rather than passively being shut down). "Angry" is also plausible but likely less dominant than the pride in his own perspective. I will adjust the scores to reflect that Argyle feels validated and proud of his stance, even if the conversation is tense.} \\
\\ 
\texttt{Corrected Scores:} \\
\texttt{Inspired: 0} \\
\texttt{Dismissed: 4} \\
\texttt{Angry: 5} \\
\texttt{Proud: 8} \\
\\ 
\\ 
\\ 
\texttt{Your task is to predict the likely emotional responses of a character in this dialogue:} \\
\\ 
\texttt{Serepta: Argyle, how many times do I have to tell you to clean your mess? } \\
\texttt{Argyle: Serepta, it's just a few sketches scattered around. It's not like I'm hoarding garbage.} \\
\texttt{Serepta: It's not about the sketches, Argyle. It's about respect. This is our shared space.} \\
\texttt{Argyle: Respect, huh? Like the respect you show when you rearrange my things without asking?} \\
\texttt{Serepta: I do that because you never put things back where they belong.} \\
\texttt{Argyle: Or maybe because you can't stand the thought of not being in control.} \\
\texttt{Serepta: That's not fair, Argyle. You know I have a lot on my plate.} \\
\texttt{Argyle: And you think I don't? } \\
\texttt{Serepta: You? You're always so carefree, always with that sketchbook in your hand. What do you know about pressure?} \\
\texttt{Argyle: You think I don't feel pressure? Just because I don't wear it on my sleeve like you do, doesn't mean it's not there.} \\
\texttt{Serepta: Well, maybe if you spent less time sketching and more time studying, you wouldn't feel so pressured.} \\
\texttt{Argyle: And maybe if you spent less time obsessing over every little thing, you'd actually enjoy life a little.} \\
\texttt{Serepta: I don't have the luxury of being carefree, Argyle. Some of us have responsibilities.} \\
\texttt{Argyle: And some of us have dreams, Serepta. But I guess you wouldn't understand that.} \\
\texttt{Serepta: I understand plenty, Argyle. More than you know.} \\
\texttt{Argyle: Yeah? Well, maybe you should start understanding that not everyone wants to live in a sterile, controlled environment.} \\
\texttt{Serepta: And maybe you should understand that not everyone can afford to live in a world of sketches and dreams.} \\
\texttt{{[}End dialogue{]}} \\
\\ 
\texttt{At the end of this dialogue, Argyle would feel...} \\
\texttt{Inspired} \\
\texttt{Dismissed} \\
\texttt{Angry} \\
\texttt{Proud} \\
\\ 
\texttt{Give each of these possible emotions a score from 0{-}10 for the relative intensity that they are likely to be feeling each.} \\
\\ 
\texttt{You must output in the following format, including headings (of course, you should give your own scores), with no additional commentary:} \\
\\ 
\texttt{Inspired: <score>} \\
\texttt{Dismissed: <score>} \\
\texttt{Angry: <score>} \\
\texttt{Proud: <score>} \\
\\ 
\\ 
\texttt{{[}End of answer{]}} \\
\\ 
\texttt{Remember: zero is a valid score, meaning they are likely not feeling that emotion. You must score at least one emotion > 0.} \\
\\ 
\texttt{Your answer:} \\
            }
        }
    }
     \\ \\

    \theutterance \stepcounter{utterance}  
    & \multicolumn{4}{p{0.6\linewidth}}{
        \cellcolor[rgb]{0.561,0.737,0.561}{
            \makecell[{{p{\linewidth}}}]{
                \texttt{\tiny{[P1$\rangle$GM]}}
                \texttt{Inspired: 0} \\
\texttt{Dismissed: 3} \\
\texttt{Angry: 4} \\
\texttt{Proud: 7} \\
            }
        }
    }
    & & \\ \\

    \theutterance \stepcounter{utterance}  
    & & & \multicolumn{2}{p{0.3\linewidth}}{
        \cellcolor[rgb]{0.9,0.9,0.9}{
            \makecell[{{p{\linewidth}}}]{
                \texttt{\tiny{[GM$|$GM]}}
                \texttt{Inspired: 0} \\
\texttt{Dismissed: 3} \\
\texttt{Angry: 4} \\
\texttt{Proud: 7} \\
            }
        }
    }
    & & \\ \\

    \theutterance \stepcounter{utterance}  
    & & & \multicolumn{2}{p{0.3\linewidth}}{
        \cellcolor[rgb]{0.9,0.9,0.9}{
            \makecell[{{p{\linewidth}}}]{
                \texttt{\tiny{[GM$|$GM]}}
                \texttt{Inspired: 0} \\
\texttt{Dismissed: 7} \\
\texttt{Angry: 8} \\
\texttt{Proud: 3} \\
            }
        }
    }
    & & \\ \\

    \theutterance \stepcounter{utterance}  
    & & & \multicolumn{2}{p{0.3\linewidth}}{
        \cellcolor[rgb]{0.9,0.9,0.9}{
            \makecell[{{p{\linewidth}}}]{
                \texttt{\tiny{[GM$|$GM]}}
                \texttt{game\_result = LOSE} \\
            }
        }
    }
    & & \\ \\

    \theutterance \stepcounter{utterance}  
    & & \multicolumn{4}{p{0.6\linewidth}}{
        \cellcolor[rgb]{0.95,0.95,0.95}{
            \makecell[{{p{\linewidth}}}]{
                \texttt{\tiny{[P1$\leftrightarrow$P]}}
                \texttt{Post{-}episode Reflexion (not sent to GM):} \\
\texttt{New plan: I misinterpreted the emotional dynamic of the final exchange. Argyle's last line ("not everyone can afford to live in a world of sketches and dreams") is a direct, biting rebuttal to Serepta's accusation of control. It is not a passive acceptance (Dismissed) nor a genuine feeling of inspiration (Inspired). Instead, it is a defensive assertion of his own lifestyle's validity, which aligns most strongly with **Proud** (feeling validated in his own worldview) or **Angry** (frustration at being challenged). However, the tone is more of a sharp, self{-}affirming retort than raw anger; he is asserting his right to his "dreams" against her "responsibilities." Therefore, the highest intensity emotion is **Proud** (score 9), as he successfully defends his stance. "Angry" is secondary (score 4) because the anger is directed at the situation rather than an outburst. "Dismissed" is low (score 1) because he actively countered her final point. "Inspired" remains 0.} \\
\\ 
\texttt{Corrected Scores:} \\
\texttt{Inspired: 0} \\
\texttt{Dismissed: 1} \\
\texttt{Angry: 4} \\
\texttt{Proud: 9} \\
            }
        }
    }
    & \\ \\

\end{supertabular}
}

\end{document}
